% -*- coding=utf-8 -*-
\chapter{附录}

\section*{一、处理碰撞事例的Python代码}
%\begin{lstlisting}[language={python}, basicstyle = \ttfamily,breaklines = true, caption={Python 源代码}, label=code_1]
\begin{lstlisting}[language=python,caption={模拟退火算法},label=Code:101]
kk=2;
[mdd,ndd]=size(dd);
while ~isempty(V)
    [tmpd,j]=min(W(i,V));tmpj=V(j);
for k=2:ndd
    [tmp1,jj]=min(dd(1,k)+W(dd(2,k),V));
    tmp2=V(jj);tt(k-1,:)=[tmp1,tmp2,jj];
end
    tmp=[tmpd,tmpj,j;tt];[tmp3,tmp4]=min(tmp(:,1));
if tmp3==tmpd, ss(1:2,kk)=[i;tmp(tmp4,2)];
 \end{lstlisting}
 
 \newpage
 
 \section*{二、处理碰撞事例的Julia代码}
\begin{lstlisting}[language=matlab,caption={模拟退火算法},label=Code:102]
using LinearAlgebra
using Random

"""
PSO求解最小特征值及特征向量
A: 对称矩阵
n_particles: 粒子数量
max_iter: 最大迭代次数
"""
function pso_min_eigen(A::Matrix{Float64}; n_particles=50, max_iter=200)
    n = size(A, 1)

    # 定义目标函数：瑞利商 f(x) = (x'Ax) / (x'x)
    function rayleigh_quotient(x)
        return (x' * A * x) / (x' * x)
    end

    # 初始化粒子位置 (随机单位向量方向) 和速度
    particles = [randn(n) for _ in 1:n_particles]
    velocities = [zeros(n) for _ in 1:n_particles]

    # 个体最优位置及其适应度
    pbest_pos = copy(particles)
    pbest_val = [rayleigh_quotient(p) for p in particles]

    # 全局最优
    gbest_idx = argmin(pbest_val)
    gbest_pos = copy(pbest_pos[gbest_idx])
    gbest_val = pbest_val[gbest_idx]

    # 超参数
    w = 0.5   # 惯性权重
    c1 = 1.5  # 个体学习因子
    c2 = 1.5  # 全局学习因子

    for _ in 1:max_iter
        for i in 1:n_particles
            # 更新速度
            velocities[i] = w * velocities[i] + 
                            c1 * rand() * (pbest_pos[i] - particles[i]) + 
                            c2 * rand() * (gbest_pos - particles[i])

            # 更新位置
            particles[i] += velocities[i]

            # 计算适应度
            val = rayleigh_quotient(particles[i])

            # 更新个体最优
            if val < pbest_val[i]
                pbest_val[i] = val
                pbest_pos[i] = copy(particles[i])

                # 更新全局最优
                if val < gbest_val
                    gbest_val = val
                    gbest_pos = copy(particles[i])
                end
            end
        end
    end

    # 返回结果：特征值与单位特征向量
    return gbest_val, gbest_pos / norm(gbest_pos)
end

# --- 测试 ---
Random.seed!(123)
A = rand(5, 5)
A = A' * A # 构造对称正定矩阵
λ, v = pso_min_eigen(A)

println("PSO计算的最小特征值: ", λ)
println("对比线性代数库计算: ", minimum(eigvals(A)))
println("单位特征向量: ", v)
\end{lstlisting}

 